\chapter{质点系动力学}

\section{质点系的动量定理和动量守恒定律}

\begin{problem}{6.1.1}
求由抛物线 $y = x^2$ 和直线 $y = a$ 围成的区域形状的均质薄板的质心.
\end{problem}
\clearpage


\begin{problem}{6.1.2}
求由四个平面: x=0, y=0, z=0 和 $4x+2y+z=8$ 为边界的均质实心体的质心。
\end{problem}
\clearpage


\begin{problem}{6.1.3}
求均质的内外半径分别为 a 和 b 的半球壳的质心.
\end{problem}
\clearpage


\begin{problem}{6.1.4}
图 6.7 所示的画有斜线的匀质薄板是由半径为 R 的大圆板挖去两个半径为 $\frac{1}{2}R$ 的完整的和不完整的小圆板形成的。两小圆的圆心 A、B 处于同一条直径上，且在圆心 O 的两侧与 O 点的距离 R 和 $\frac{R}{2}$ 处。求此薄板的质心。
\begin{center}
\begin{tikzpicture}[>=Stealth, thick]
  \def\R{2.1}

  % Hatched plate inside large circle, outside the two holes
  \begin{scope}
    \clip (0,0) circle (\R);
    \foreach \k in {-14,-13,...,14} {
      \draw[thin] (-2.6, \k*0.26 - 0.9) -- (2.6, \k*0.26 + 0.9);
    }
    \fill[white] (0,\R) circle (0.5*\R);
    \fill[white] (0,-0.5*\R) circle (0.5*\R);
  \end{scope}

  % Circle boundaries
  \draw (0,0) circle (\R);
  \begin{scope}
    \clip (0,0) circle (\R);
    \draw (0,\R) circle (0.5*\R);
  \end{scope}
  \draw (0,-0.5*\R) circle (0.5*\R);

  % Coordinate axes
  \draw[->] (0,0) -- (2.6,0) node[below] {$x$};
  \draw[->] (0,0) -- (0,2.7) node[left] {$y$};
  \node[below right, inner sep=1pt] at (0.08,-0.08) {$O$};

  % Labels A, B and radius arrows
  \node[above right=1pt] at (0.04,\R) {$A$};
  \draw[->] (0,\R) -- ++(-60:0.5*\R) node[pos=0.60, right=2pt, inner sep=1pt] {$\frac{1}{2}R$};

  \draw[->] (0,0) -- (22:\R) node[pos=0.65, above, fill=white, inner sep=1.5pt] {$R$};

  \node[left=2pt] at (-0.06,-0.5*\R) {$B$};
  \draw[->] (0,-0.5*\R) -- ++(-18:0.5*\R) node[pos=0.58, above=2pt, inner sep=1pt] {$\frac{1}{2}R$};
\end{tikzpicture}

\end{center}
\end{problem}
\clearpage


\begin{problem}{6.1.5}
正圆锥体底面的半径为 $a$ ，今用一通过其对称轴的平面将它等分为二。证明任一半的质心和该轴的距离为 $\frac{a}{\pi}$ 。
\begin{center}
\begin{tikzpicture}[>=Stealth, thick]
  % Coordinates
  \coordinate (O) at (0,0);
  \coordinate (C) at (0,-3.2);     % Base center
  \coordinate (L) at (-1.75,-3.2); % Left base point
  \coordinate (R) at (1.75,-3.2);  % Right base point
  % Split plane intersection with base ellipse (x_rad=1.75, y_rad=0.52) at -65 deg / +115 deg
  \coordinate (BF) at ({1.75*cos(-65)}, {-3.2 + 0.52*sin(-65)});
  \coordinate (BB) at ({1.75*cos(115)}, {-3.2 + 0.52*sin(115)});

  % Upper cross-section at ratio t = 0.42 from O
  \def\t{0.42}
  \coordinate (Cz) at (0, -3.2*\t);
  \coordinate (Lz) at (-1.75*\t, -3.2*\t);
  \coordinate (Rz) at (1.75*\t, -3.2*\t);
  \coordinate (BFz) at ({1.75*\t*cos(-65)}, {-3.2*\t + 0.52*\t*sin(-65)});
  \coordinate (BBz) at ({1.75*\t*cos(115)}, {-3.2*\t + 0.52*\t*sin(115)});

  % Axes at O
  \draw[->] (O) -- (-2.1,0) node[left] {$y$};
  \draw[->] (O) -- (1.6,-1.15) node[below right] {$x$};
  \draw[->] (O) -- (0,-4.4) node[right] {$z$};
  \node[above right] at (O) {$O$};

  % Cone outer generators
  \draw (O) -- (L);
  \draw (O) -- (BF);
  \draw[dashed] (O) -- (BB);

  % Base ellipse: front/right solid, back-left dashed
  \draw (L) arc[start angle=180, end angle=360, x radius=1.75, y radius=0.52]
        arc[start angle=0, end angle=115, x radius=1.75, y radius=0.52];
  \draw[dashed] (L) arc[start angle=180, end angle=115, x radius=1.75, y radius=0.52];

  % Base radii and labels
  \draw (L) -- (C) node[midway, below] {$a$};
  \draw[dashed] (BB) -- (BF);
  \node[right] at (0.04,-3.12) {$h$};

  % Cross-section at height z
  \draw (Lz) arc[start angle=180, end angle=295, x radius=1.75*\t, y radius=0.52*\t];
  \draw[dashed] (Lz) arc[start angle=180, end angle=115, x radius=1.75*\t, y radius=0.52*\t];
  \draw[dashed] (BFz) arc[start angle=-65, end angle=115, x radius=1.75*\t, y radius=0.52*\t];
  \draw[dashed] (Lz) -- (Cz);
  \draw[dashed] (BBz) -- (BFz);

  % Label for cross-section radius
  \draw[thin] (-0.32, -3.2*\t + 0.06) -- (-1.05, -3.2*\t + 0.06) node[left] {$\dfrac{a}{h}z$};
\end{tikzpicture}

\end{center}
\end{problem}
\clearpage


\begin{problem}{6.1.6}
证明一个均质的三角形薄板的质心与位于该三角形三个顶点有相等质量的质点系统的质心重合.
\end{problem}
\clearpage


\begin{problem}{6.1.7}
修理街道电网用的迁移式转动起重机安装在质量为 $1\mathrm{t}$ 的自动车上，质量 $200\mathrm{kg}$ 的起重机的摇架 $K$ 装在质量 $100\mathrm{kg}$ 、杆长 $4\mathrm{m}$ 的杆 $L$ 上. 可绕与图面垂直的水平轴 $O$ 转动. 摇架的质心 $C$ 离轴 $O$ 的距离 $OC = 3.5\mathrm{m}$ .

(1) 开始起重机杆 L 处于水平位置, 自动车处于静止. 求起重机杆转过 $60^{\circ}$ 时, 未上闸

的自动车移动的距离,运动时所受阻力均可不计;

(2) 自动车上了闸, 起重机杆从水平位置以恒定的角速度 $\omega = 0.02\mathrm{rad}\cdot \mathrm{s}^{-1}$ 转到 $60^{\circ}$ 期间, 靠地面的摩擦力, 自动车保持不动. 求在此期间摩擦力给予的冲量;

(3) 若摩擦因数 $\mu=0.1$ ，要在转动过程中保持自动车不动，对转动角速度 $\omega$ 有何限制？
\begin{center}
\begin{tikzpicture}[>=Stealth, thick]
  % Ground with hatching
  \draw (-0.3,0) -- (5.2,0);
  \foreach \x in {-0.15,0.10,...,5.10} {
    \draw[thin] (\x,0) -- (\x-0.16,-0.16);
  }

  % Truck body
  % Flatbed, cab, and hood outline
  \draw (0.05,0.52) -- (0.05,1.12) -- (2.75,1.12)
        -- (2.82,1.68) .. controls (2.85,1.88) and (3.00,1.92) .. (3.25,1.92)
        -- (3.55,1.92) .. controls (3.72,1.92) and (3.80,1.75) .. (3.86,1.45)
        -- (4.55,1.42) .. controls (4.78,1.40) and (4.85,1.22) .. (4.86,0.92)
        -- (4.96,0.92) -- (4.96,0.72) -- (4.86,0.72) -- (4.86,0.52) -- cycle;

  % Cab window and door details
  \draw (3.12,1.12) -- (3.12,1.66) .. controls (3.12,1.76) and (3.22,1.80) .. (3.38,1.80)
        -- (3.54,1.80) .. controls (3.64,1.80) and (3.70,1.68) .. (3.70,1.52)
        -- (3.70,1.12);
  \draw (3.12,1.44) -- (3.70,1.44);

  % Wheel arch above front wheel
  \fill[white] (4.25,0.42) ++(162:0.53) arc[start angle=162, end angle=18, radius=0.53] -- cycle;
  \draw (4.25,0.42) ++(165:0.52) arc[start angle=165, end angle=15, radius=0.52];

  % Wheels (white fill covers chassis bottom line cleanly)
  \foreach \wx in {0.88,1.98,4.25} {
    \fill[white] (\wx,0.42) circle (0.42);
    \draw (\wx,0.42) circle (0.42);
    \draw (\wx,0.42) circle (0.17);
  }

  % Crane pedestal on flatbed
  \draw (0.55,1.12) rectangle (1.62,1.34);
  \draw (0.38,1.34) rectangle (1.78,1.44);
  \draw (0.56,1.44) -- (0.56,2.04) (0.68,1.44) -- (0.68,2.04);

  % Hydraulic cylinder
  \draw (0.68,2.05) -- (1.15,1.72);
  \draw (1.02,1.85) -- (1.42,1.57) -- (1.35,1.48) -- (0.95,1.76) -- cycle;

  % Pivot O and boom L (angle 18 deg)
  \coordinate (O) at (0.62,2.10);
  \def\ang{18}
  \coordinate (C) at ($(O)+(\ang:3.05)$);
  \coordinate (E) at ($(O)+(\ang:3.36)$);

  % Boom double lines
  \draw ($(O)+(\ang+90:0.045)$) -- ($(O)+(\ang:2.74)+(\ang+90:0.045)$);
  \draw ($(O)+(\ang-90:0.045)$) -- ($(O)+(\ang:2.74)+(\ang-90:0.045)$);
  \node[below] at ($(O)+(\ang:1.55)+(\ang-90:0.12)$) {$L$};

  % Pivot circle
  \fill[white] (O) circle (0.08);
  \draw (O) circle (0.08);
  \node[below left] at (O) {$O$};

  % Cradle K and center of mass C
  \begin{scope}[shift={(C)}, rotate=\ang]
    \draw (-0.31,-0.19) rectangle (0.31,0.19);
    \draw (0,0) circle (0.065);
    \draw (0.12,0.065) -- (0.24,0.065) (0.12,-0.065) -- (0.24,-0.065);
    \node[below=6pt] at (0,-0.16) {$C$};
    \node[right=4pt] at (0.31,0) {$K$};
  \end{scope}

  % Dimension lines (3.5 m to C, 4 m to end E)
  \coordinate (O1) at ($(O)+(\ang+90:0.55)$);
  \coordinate (O2) at ($(O)+(\ang+90:0.95)$);
  \coordinate (C2) at ($(C)+(\ang+90:0.95)$);
  \coordinate (E1) at ($(E)+(\ang+90:0.55)$);

  \draw[thin] ($(O)+(\ang+90:0.16)$) -- ($(O)+(\ang+90:1.08)$);
  \draw[thin] ($(C)+(\ang+90:0.22)$) -- ($(C)+(\ang+90:1.08)$);
  \draw[thin] ($(E)+(\ang+90:0.22)$) -- ($(E)+(\ang+90:0.68)$);

  \draw[<->, thin] (O2) -- (C2) node[midway, above, rotate=\ang] {$3.5\mathrm{m}$};
  \draw[<->, thin] (O1) -- (E1) node[midway, above, rotate=\ang] {$4\mathrm{m}$};
\end{tikzpicture}

\end{center}
\end{problem}
\clearpage


\begin{problem}{6.1.8}
一火箭垂直向上发射, 当它达到最高点时炸裂成三个等质量的碎片. 观察到其中一块碎片经时间 $t_1$ 垂直地落到地上, 而其他两块碎片在炸裂后的 $t_2$ 时刻落到地上, 求炸裂时离地面的高度.
\end{problem}
\clearpage


\begin{problem}{6.1.9}
以下各量均用万有引力常数 $G$ 、距离 $L$ 和质量表示：

(1) 相距为 L 的某一等质量 $(M_{1}=M_{2}=M)$ 双星的转动周期；

(2) 相距为 L 的某一不等质量 $(M_{1} \neq M_{2})$ 双星的转动周期；

(3) 边长为 L 的某一等质量等边三角形三星的转动周期；

(4) 某一不等质量 $(M_{1} \neq M_{2} \neq M_{3})$ 的等边 (边长为 $L$ ) 三角形三星的转动周期.
\end{problem}
\clearpage


\begin{problem}{6.1.10}
一质量为 $M$ 的有轨平板车能无摩擦地在一条水平直轨道上运动. 初始时刻, $N$ 个质量均为 $m$ 的人站在静止的平板车上.

(1) 这 N 个人一起跑向车的一头, 并且同时跳下车, 在跳下车前的瞬时, 他们相对于车子的速度为 $v_{r}$ , 求这 N 个人跳下后平板车的速度;

(2) 如果这 $N$ 个人是一个接一个地跑离车子 (每次仅有一个在跑). 每个人在跳下车

前相对于车子的速度都为 $v_{r}$ . 求这车子的最终速度；

(3) 在上述两种情况下,哪种情况下车子达到的速度大?
\end{problem}
\clearpage


\begin{problem}{6.1.11}
证明质量为 $M$ 的飞机在水平航线上飞行时落下一个质量为 $m$ 的炸弹时，飞机将有一个向上的加速度 $mg / M$ .
\end{problem}
\clearpage


\begin{problem}{6.1.12}
水从 5m 高处以每分钟 60kg 的速率注入静放在地上的桶中, 桶的质量为 12kg. 求水注入一分钟时桶受到地面的作用力.
\end{problem}
\clearpage


\begin{problem}{6.1.13}
质量为 $m_A$ 与 $m_B$ 的两个小球 $A$ 和 $B$ , 用不可伸长的轻绳相连, 静止在水平面上. 绳子拉紧时, 给 $B$ 球一个冲量 $I, I$ 的方向与 $AB$ 成 $\alpha$ 角, $\alpha < \frac{\pi}{2}$ , 如图 6.13 所示. 已知冲击后 $B$ 的速度 $\bm{v}_B$ 之大小, 求 $\bm{v}_B$ 的方向和冲量 $I$ 的大小.
\begin{center}
\begin{tikzpicture}[>=Stealth, thick]
  \coordinate (A) at (0,0);
  \coordinate (B) at (2.8,0);
  \coordinate (C) at (4.1,0);
  \coordinate (Iend) at ($(B)+(36:1.45)$);

  \draw (A) -- (B);
  \draw[dashed] (B) -- (C);
  \draw[->] (B) -- (Iend) node[above] {$I$};

  \fill (A) circle (2.2pt);
  \fill (B) circle (2.2pt);
  \node[below=3pt] at (A) {$A$};
  \node[below=3pt] at (B) {$B$};

  \pic[draw, thin, angle radius=0.55cm, angle eccentricity=1.35, "$\alpha$"] {angle=C--B--Iend};
\end{tikzpicture}

\end{center}
\end{problem}
\clearpage


\begin{problem}{6.1.14}
扫雪机在水平路上以速度 15km/h 行驶, 每分钟把质量为 30t 的雪抛到路旁, 雪相对于雪铲的速度为 12m/s, 并沿 45°角方向. 求路面与轮胎间的横向力, 并求驱动扫雪机工件所需之牵引力.
\end{problem}
\clearpage


\begin{problem}{6.1.15}
有 $m_{1}, m_{2}$ 和 $m_{3}$ 三质量受到它们的相互间万有引力的作用, 如果最初这些质量位于边长为 $a$ 的等边三角形的顶点, 只要这些质量的初速是适当选择的, 它们的位置就会连续确定一个边长为 $a$ 以等角速 $\omega$ 转动的等边三角形. 求 $\omega$ 与 $m_{1}, m_{2}, m_{3}, a$ 以及引力常数 $G$ 的关系.
\end{problem}
\clearpage


\begin{problem}{6.1.16}
一个仅由自身引力保持在一起的小天体. 如果它非常靠近大质量体的话, 可以被这个大天体引起的潮汐力所分裂. 求一个直径为 $1\mathrm{km}$ 、密度为 $\rho = 2\mathrm{g/cm}^3$ 的小天体围绕地球做圆周运动而不被拉散的圆周极限半径. 已知地球质量 $M = 5.97 \times 10^{27}\mathrm{g}$ .
\end{problem}
\clearpage


\section{质点系的角动量定理和角动量守恒定律}

\begin{problem}{6.2.1}
半径为 R、质量为 M 的水平均质圆盘可绕通过其中心的铅垂轴无摩擦地转动.

(1) 若圆盘绕该轴以 $\Omega$ 的角速度转动, 求圆盘对该轴的角动量;

(2) 质量为 $m$ 的甲虫按 $s = \frac{1}{2} at^2 (a$ 为常量) 的规律沿圆盘的边缘爬行, 开始时两者都是静止的, 求甲虫爬行后圆盘的角速度;

(3) 质量为 $m$ 的甲虫从边缘上一点突然以对圆盘恒定的角速度 $\omega$ 作半径为 $\frac{1}{2} R$ 的圆周运动, 开始两者都是静止的. 求甲虫突然爬行后圆盘的角速度.
\end{problem}
\clearpage


\begin{problem}{6.2.2}
一个半径为 $a$ 的均质圆盘可在竖直平面内围绕通过它的中心 $C$ 的水平轴自由转动, 质量为圆盘 $\frac{1}{10}$ 的昆虫 $A$ 在盘的最低点, 整个系统处于静止状态. 突然, 昆虫以相对于圆盘不变的速率 $v$ 开始沿边缘爬行, 证明圆盘的初角速度为 $v / 6a$ ; 如用 $\theta$ 表示 $CA$ 线与竖直向下方向间的夹角, 证明微分方程 $6a \ddot{\theta} + g \sin \theta = 0$ 成立, 证明圆盘对昆虫的作用

力在垂直于半径方向的分量为 $\frac{1}{12}mg\sin\theta$ ，其中m是圆盘的质量.
\begin{center}
\begin{tikzpicture}[>=Stealth, thick]
  \def\a{1.95}
  \coordinate (C) at (0,0);
  \coordinate (V) at (0,-\a-0.45);
  \coordinate (P) at (-112:\a);
  \coordinate (A) at (-48:\a);

  % Disk
  \draw (C) circle (\a);
  \node[above=2pt] at (C) {$C$};

  % Radial and vertical lines
  \draw (C) -- (V);
  \draw (C) -- (P);
  \draw (C) -- (A) node[midway, above right] {$a$};
  \node[right=3pt] at (A) {$A$};

  % Angles phi and theta
  \pic[draw=none, angle radius=0.95cm, angle eccentricity=0.82, "$\varphi$"] {angle=P--C--V};
  \pic[draw=none, angle radius=0.95cm, angle eccentricity=0.80, "$\theta$"] {angle=V--C--A};

  % Forces at A: weight (1/10 mg) and tangential force F
  \draw[->] (A) -- ++(0,-1.15) node[pos=0.8, right=2pt] {$\dfrac{1}{10}mg$};
  \draw[->] (A) -- ++(42:0.95) node[above right=-1pt] {$F$};
\end{tikzpicture}

\end{center}
\end{problem}
\clearpage


\begin{problem}{6.2.3}
质量为 $M$ 、半径为 $a$ 的均质圆柱可围绕它在水平方向固定的对称轴自由转

动, 质量为 $m$ 、长度 $b = 2\pi an (n$ 为整数) 的均质链条一端连在圆柱上, 然后绕圆柱 $n$ 圈, 让非常小的一段链条自由挂着, 整个系统由静止释放. 求圆柱的角速度与圆柱角位移 $\theta$ 的函数关系.
\end{problem}
\clearpage


\begin{problem}{6.2.4}
质量均为 $m$ 的两个质点用一根长度为 $a$ 的刚性的轻杆联结, 开始系统静止在一光滑水平面上. 在 $t = 0$ 时, 一力 $\bm{F}$ 作用于其中一个质点, 大小保持不变, 方向保持与杆垂直, 用角动量定理和动能定理两种方法求杆与开始方向夹角为 $\varphi$ 时的角速度.
\end{problem}
\clearpage


\begin{problem}{6.2.5}
质量分别为 $m$ 和 $M$ 的质点 $A$ 和 $B$ , 如质点 $A$ 相对于质点 $B$ 的角动量为 $J_{AB}$ , 证明: 在质心平动参考系中, 系统相对于质心的角动量为 $\frac{M}{M + m} J_{AB}$ .
\end{problem}
\clearpage


\begin{problem}{6.2.6}
在赤道无风带有一艘因无风而停着的船. 船长决定将 $m=200kg$ 的锚升到长 l=20m 的桅杆顶上的权宜之计使船运动, 船及其负载(不包括锚)的质量 $M=1000kg$ , 问:

(1) 船为什么能开始运动?

(2) 说明船向哪个方向运动?

(3) 锚升到桅杆顶后船的运动速度多大?
\end{problem}
\clearpage


\begin{problem}{6.2.7}
分别考虑地球-月球系统和火星-火星的一个卫星系统,为简单起见,考虑每个系统时不考虑它们与其他星体之间的相互作用.月球绕地球的运动比地球的自转慢得多,而火星的一个卫星绕火星的运动比火星的自转快,哪种情况卫星引起行星上的潮汐使卫星、行星之间的距离增加,而另一种情况距离减少?
\end{problem}
\clearpage


\section{质点系的动能定理和机械能守恒定律}

\begin{problem}{6.3.1}
一系统由两个轮子、杆 $O_{1}O_{2}$ 及其平行杆 $AB$ 组成，轮子半径为 $r$ 、质量为 $M_{1}$ ，均匀分布于边缘， $O_{1}O_{2}$ 连接于两轮轴上， $AB$ 杆与固连于轮子的 $O_{1}A, O_{2}B$ 铰接， $O_{1}O_{2}$ 杆及 $AB$ 杆质量均为 $M_{2}, O_{1}A, O_{2}B$ 长为 $\frac{1}{2} r$ ，质量可以不计。若在图6.21所示位置时，轮子在直线轨道上做纯滚动，轮轴的速度为 $v_{0}$ ，求此系统的动能。
\begin{center}
\begin{tikzpicture}[>=Stealth, thick]
  % Ground with hatching
  \draw (-1.3,0) -- (4.5,0);
  \foreach \x in {-1.15,-0.90,...,4.40} {
    \draw[thin] (\x,0) -- (\x-0.16,-0.16);
  }

  \def\r{1.05}
  \coordinate (O1) at (0,\r);
  \coordinate (O2) at (3.2,\r);
  \coordinate (A) at ($(O1)+(45:0.55*\r)$);
  \coordinate (B) at ($(O2)+(45:0.55*\r)$);
  \coordinate (H) at ($(O2)+(0.85,0)$);

  % Two wheels
  \draw (O1) circle (\r);
  \draw (O2) circle (\r);

  % Radius r arrow on left wheel
  \draw[->] (O1) -- ++(155:\r) node[midway, below left=1pt] {$r$};

  % Linkage rods O1-O2, O1-A, O2-B, A-B
  \draw (O1) -- (O2);
  \draw (O1) -- (A) node[pos=0.52, above left=-1pt, inner sep=1pt] {$\frac{1}{2}r$};
  \draw (O2) -- (B);
  \draw (A) -- (B);

  % Dashed horizontal extension at O2 and angle phi
  \draw[dashed] (O2) -- (H);
  \pic[draw, thin, angle radius=0.45cm, angle eccentricity=1.42, "$\varphi$"] {angle=H--O2--B};

  % Vertex labels
  \node[below=2pt] at (O1) {$O_1$};
  \node[below=2pt] at (O2) {$O_2$};
  \node[above=2pt] at (A) {$A$};
  \node[above=2pt] at (B) {$B$};
\end{tikzpicture}

\end{center}
\end{problem}
\clearpage


\begin{problem}{6.3.2}
质量均为 $m$ 的两质点被一原长为 $l$ 、劲度系数为 $k$ 的弹簧连接，原在光滑的水平面上处于静止，在 $t = 0$ 时，一个质点受到一个垂直于弹簧方向的冲量 $I$ ，假定在以后的运动中弹簧始终沿着两质点的连线（即不弯曲）.

(1) t 时刻两质点的总能量与总动量多大?

(2) 质心速度和相对于质心的总角动量为何值?

（3）若 $m = 1\mathrm{kg}, k = 1\mathrm{N / m}, l = 1\mathrm{m}, I = 1\mathrm{N} \cdot \mathrm{s}$ ，冲量作用后的运动过程中，两质点之间的最大距离多大？

(4) 两质点的最大瞬时速率多大?
\end{problem}
\clearpage


\begin{problem}{6.3.3}
一质量为 $m$ 的卫星围绕地球以速度 $\bm{v}$ 在一半径为 $R$ 的圆轨道上运动. 卫星突然吸附了一小质量 $\delta m$ , 这块小质量在发生碰撞前是静止不动的. 求此卫星总能量的改变量, 并求这个新轨道的长半轴.
\end{problem}
\clearpage


\begin{problem}{6.3.4}
两根质量为 $m$ 、长 $2a$ 完全相同的均质杆 $AB$ 和 $BC$ 于 $B$ 点自由铰接，在一光滑的水平桌上运动。用两杆的夹角 $2\theta$ 、 $AC$ 和平面上固定方向间的夹角 $\varphi$ 及其对时间的导

数表示系统在其质心平动参考系中的动能.
\begin{center}
\begin{tikzpicture}[>=Stealth, thick]
  % Geometry: rods BA and BC of length 2a, angle ABC = 2*theta
  \coordinate (B) at (-2.0,-0.95);
  \coordinate (A) at ($(B)+(44:3.6)$);
  \coordinate (C) at ($(B)+(-5:3.6)$);
  \coordinate (D) at ($(B)!0.5!(A)$);
  \coordinate (E) at ($(B)!0.5!(C)$);
  \coordinate (M) at ($(A)!0.5!(C)$);
  \coordinate (O) at ($(D)!0.5!(E)$);

  % Rods AB and BC
  \draw (A) -- (B) -- (C);
  \node[above right] at (A) {$A$};
  \node[below left] at (B) {$B$};
  \node[below right] at (C) {$C$};
  \node[above left] at (D) {$D$};
  \node[below=2pt] at (E) {$E$};
  \node[below=3pt] at ($(B)!0.5!(E)$) {$a$};

  % Dashed median BM and segments DE, AC
  \draw[dashed] (A) -- (C);
  \draw[dashed] (D) -- (E);
  \draw[dashed] (B) -- (M);

  % Center-of-mass axes at O
  \coordinate (Xtip) at ($(O)+(2.45,0)$);
  \coordinate (Ytip) at ($(O)+(0,1.75)$);
  \draw[->] (O) -- (Xtip) node[below] {$x$};
  \draw[->] (O) -- (Ytip) node[left] {$y$};
  \fill (O) circle (1.4pt);
  \node[below left=2pt] at (O) {$O$};

  % Angle theta at B
  \pic[draw=none, angle radius=0.85cm, angle eccentricity=0.85, "$\theta$"] {angle=O--B--A};

  % Intersection of x-axis and AC for angle phi
  \coordinate (P) at (intersection of O--Xtip and A--C);
  \pic[draw, thin, angle radius=0.38cm, angle eccentricity=1.45, "$\varphi$"] {angle=Xtip--P--A};
\end{tikzpicture}

\end{center}
\end{problem}
\clearpage


\begin{problem}{6.3.5}
两个质量均为 $m$ 、速率均为 $u$ 的粒子做斜碰，如图6.24所示.

(1) 求质心平动参考系的速度V；

(2) 将在质心平动参考系中的总能量与在静参考系中的总能量进行比较.
\begin{center}
\begin{tikzpicture}[>=Stealth, thick]
  \coordinate (O) at (0,0);
  \coordinate (M1) at (-2.15,0);
  \coordinate (U1) at (-1.15,0);
  \coordinate (Xtip) at (2.4,0);
  \coordinate (Ytip) at (0,1.55);
  \coordinate (M2) at (24:2.25);
  \coordinate (U2) at (24:1.30);

  % Axes
  \draw[->] (O) -- (Xtip) node[below] {$x$};
  \draw[->] (O) -- (Ytip) node[left] {$y$};
  \node[below=2pt] at (O) {$O$};

  % Particle m1 and velocity u
  \fill (M1) circle (2.8pt) node[above=3pt] {$m_1$};
  \draw[->] (M1) -- (U1) node[midway, above=2pt] {$u$};
  \draw[dashed] (U1) -- (O);

  % Particle m2 and velocity u
  \fill (M2) circle (2.8pt) node[right=3pt] {$m_2$};
  \draw[->] (M2) -- (U2) node[midway, above left=1pt] {$u$};
  \draw[dashed] (U2) -- (O);

  % Angle theta_0
  \pic[draw=none, angle radius=1.05cm, angle eccentricity=0.92, "$\theta_0$"] {angle=Xtip--O--M2};
\end{tikzpicture}

\end{center}
\end{problem}
\clearpage


\begin{problem}{6.3.6}
一个质量为 $M$ 、半径为 $R$ 的匀质圆柱体置于光滑的水平面上，可绕其竖直的对称轴自由转动，两个质量均为 $m$ 的质点，用长度均为 $a$ 的不可伸长的轻绳连接在圆柱底部的一条直径两端的钩子上，将绳子按顺时针方向完全绕在圆柱体上，再用钩子将两质点直接挂在圆柱体上。证明：整个系统不论以什么角速度绕对称轴逆时针转动时，先让直接钩住质点的两钩子脱钩，只要绳子长度合适，当绳子完全展开时再让直径两端的钩子脱钩，圆柱体都将停止转动。求绳长 $a$ 与 $m$ 、 $M$ 、 $R$ 的关系。
\end{problem}
\clearpage


\begin{problem}{6.3.7}
在光滑的水平桌上,有两个质量为 $m$ 的质点,用一根长度为 $a$ 的不可伸长的轻绳相连结.两质点原先静止,绳子处于拉直状态, $t=0$ 时,给一个质点垂直于绳的冲量 $I$ .求证此后两质点做旋轮线运动,该质点与另一质点的能量之比为 $\cot^{2}\left(\frac{It}{2ma}\right)$ .
\end{problem}
\clearpage


\begin{problem}{6.3.8}
两个质量均为 $m$ 的质点，连在一根长 $l$ 、质量可忽略不计的刚性杆的两端，该系统在光滑水平面上以质心速度为 $\bm{v}$ 、角速度为 $\omega$ （其方向垂直于运动平面）运动时，其一端的一个质点与另一个质量也为 $m$ 的静止质点做完全非弹性碰撞，碰撞时运动的质点的速度方向与质心的速度方向相反。求：

(1) 碰撞以后的角速度；

(2) 在碰撞中损失的机械能.
\end{problem}
\clearpage


\begin{problem}{6.3.9}
质量均为 $m$ 的两个质点用长度为 $4a$ 的不可伸长、不计质量的绳子连结，开始两质点均以与连线相垂直的速度 $v$ 在光滑的水平面上运动，两质点轨道间距离为 $4a$ . 突然遇到一个直径可忽略的光滑钉子，此时钉子距质点1的距离为 $a$ . 若以后两质点在相反方向转动时不相碰，问在以后的运动中：

(1) 质点 1 离钉子的最大距离多大?

(2) 质点 1 离钉子最大距离时绳子张力多大?
\end{problem}
\clearpage


\begin{problem}{6.3.10}
质量为 $M$ 、半径为 $a$ 的半球, 其底面放在光滑的水平面上. 有一质量为 $m$ 的质点沿此半球下滑, 初始位置与球心的连线和铅垂线成 $\alpha$ 角, 系统开始是静止的, 求质点绕球心的角速度.
\begin{center}
\begin{tikzpicture}[>=Stealth, thick]
  \def\a{1.75}
  \coordinate (O) at (0,0);
  \coordinate (C) at (-0.85,0);
  \coordinate (Xleft) at (-2.95,0);
  \coordinate (Xright) at (2.15,0);
  \coordinate (Ytop) at (0,2.25);

  % Axes
  \draw[->] (Xleft) -- (Xright);
  \node[below] at (-2.85,0) {$x$};
  \draw[->] (O) -- (Ytop);
  \node[below=2pt] at (O) {$O$};

  % Displaced hemisphere (solid) centered at C
  \draw (C) ++(\a,0) arc[start angle=0, end angle=180, radius=\a];
  \coordinate (Ctop) at ($(C)+(0,\a+0.45)$);
  \coordinate (Ptheta) at ($(C)+(20:\a)$);
  \draw[dashed] (C) -- (Ctop);
  \draw[dashed] (C) -- (Ptheta);
  \pic[draw, thin, angle radius=0.45cm, angle eccentricity=1.45, "$\theta$"] {angle=Ptheta--C--Ctop};

  % Initial hemisphere (dashed) centered at O
  \draw[dashed] (O) ++(\a,0) arc[start angle=0, end angle=180, radius=\a];
  \coordinate (Palpha) at ($(O)+(62:\a)$);
  \draw[dashed] (O) -- (Palpha);
  \pic[draw=none, angle radius=0.82cm, angle eccentricity=0.92, "$\alpha$"] {angle=Palpha--O--Ytop};
\end{tikzpicture}

\end{center}
\end{problem}
\clearpage


\begin{problem}{6.3.11}
一根长为 $2a$ 不可伸长的轻绳，中点连着一个质量为 $M$ 的质点，两端各连一个质量为 $m$ 的质点，在一光滑的水平面上成一正好拉直的直线。给 $M$ 一个垂直于绳的水平初速 $v$ 。证明：若时间 $\tau$ 以后两端的质点发生碰撞时，质量为 $M$ 的质点从开始位置位移 $b = \frac{Mv\tau + 2ma}{M + 2m}$ ，并证明正要发生碰撞之前，绳子张力 $T = \frac{mM^2v^2}{(M + 2m)^2a}$ 。
\end{problem}
\clearpage


\begin{problem}{6.3.12}
一质量为 M、半径为 R 的重星以速度 V 通过质量密度为 $\rho$ 的非常稀薄的气体，由于它的引力场，重星将吸引迎面接近它的气体分子，并将它俘获在表面上，相对于 V，气体分子的热运动可以忽略，气体分子间的相互作用力也可以忽略不计。求星体在此刻受到来自气体分子的阻力。
\end{problem}
\clearpage


\section{两体问题}

\begin{problem}{6.4.1}
两颗星的质量分别为 $M$ 和 $m$ , 相距为 $d$ , 两星绕着不动的质心做圆周运动. 两颗星近似为质点, 质量为 $M$ 的星发生爆炸, 损失质量 $\Delta M$ , 爆炸是瞬时的、球对称的, 对残余体的合力为零, 爆炸时对另一颗星无直接作用. 证明: 当 $\Delta M < \frac{1}{2} (M + m)$ 时, 余下的双星系统的运动仍是有界的.
\end{problem}
\clearpage


\begin{problem}{6.4.2}
有两个质量均为 $m$ 的质点用一根劲度系数为 $k$ 、未伸长时长 $a$ 的轻弹簧连起来, 静置在光滑的水平面上. 第三个质量也为 $m$ 的质点以垂直于弹簧的速度 $V$ 运动, 并射中一个质点粘在上面. 要使弹簧拉伸到最大长度为 $3a, V$ 应多大?
\end{problem}
\clearpage


\begin{problem}{6.4.3}
一个质量为 $m$ 的中子以速率 $v$ 与一个质量为 $M$ 的静止的原子发生弹性碰撞. 试用静参考系和作为两体问题引入折合质量的办法, 证明碰后原子的最大速率为 $V = \frac{2m}{m + M} v$ .
\end{problem}
\clearpage


\begin{problem}{6.4.4}
两个质量分别为 $m_{1}$ 和 $m_{2}$ 的行星围绕太阳做长半轴分别为 $a_{1}$ 和 $a_{2}$ 的椭圆轨道运动, 证明它们的周期比为

\[
\frac {T _ {1}}{T _ {2}} = \left[ \frac {a _ {1} ^ {3} (M + m _ {2})}{a _ {2} ^ {3} (M + m _ {1})} \right] ^ {1 / 2}
\]

其中 M 是太阳的质量.
\end{problem}
\clearpage


\begin{problem}{6.4.5}
质量分别为 $4m$ 和 $m$ 的两质点 $A$ 和 $B$ 能沿 $x$ 轴无摩擦地运动, 两质点间有大小为 $kr$ 的相互吸引力, 其中 $k$ 是正值常量, $r$ 是两质点间的距离. $t = 0$ 时, 质点 $A$ 位于 $x = 5a$ 处, 质点 $B$ 位于 $10a$ 处, 均处于静止状态. 求:

(1) 两质点发生碰撞的 x 值；

(2) 在将发生碰撞的时刻两质点的相对速度；

(3) 如碰撞是弹性的, 碰撞后到再次发生碰撞经历的时间.
\end{problem}
\clearpage


\begin{problem}{6.4.6}
质量分别为 $m$ 和 $M$ 的两个质点分别连在一个原长为 $a$ 、劲度系数为 $k$ 的轻弹簧的两端，开始两质点均处于静止状态， $m$ 位于 $M$ 正上方 $a$ 处。 $t = 0$ 时给 $m$ 一个竖直向上的速率 $v$ ，求以后任何时刻两质点的位置。
\end{problem}
\clearpage


\begin{problem}{6.4.7}
质量为 $m$ 、带电荷 $e$ 的两个相同的带电粒子，开始相距很远，一个粒子是静止的，另一个以速率 $v$ 、瞄准距离 $\rho = 2ke^2 / mv^2$ 向静止的粒子方向运动。求质心系中和实验室参考系中的散射角 $\theta_c, \theta_r$ ，两粒子的最近距离及最近距离时两粒子的相对速率和绝对速率，其中 $k$ 是 $F = \frac{ke^2}{r^2}$ 中的系数。
\end{problem}
\clearpage


\section{变质量质点的运动}

\begin{problem}{6.5.1}
初始质量为 $m_{0}$ 的火箭在均匀重力场中以不变加速度 ng (其中 g 是重力加速度) 铅直地上升, 不计空气阻力, 燃气喷射的相对速率不变, 求火箭质量的变化规律.
\end{problem}
\clearpage


\begin{problem}{6.5.2}
火箭在均匀重力场中以不变的加速度 a 向上运动, 不计空气阻力, 燃气喷射的相对速率 $v_{r}$ 不变. 求火箭质量降为二分之一所需的时间.
\end{problem}
\clearpage


\begin{problem}{6.5.3}
初速为零的变质量质点,以不变的加速度 a 沿水平方向运动,燃气喷射的相对速率 $v_{r}$ 是常量,不计阻力,求质量减为 $\frac{1}{k}$ 时所走过的路程.
\end{problem}
\clearpage


\begin{problem}{6.5.4}
如上题的变质量质点受到滑动摩擦力, 摩擦因数为 $\mu$ . 求所走过的路程.
\end{problem}
\clearpage


\begin{problem}{6.5.5}
初始质量为 $m_{0}$ 的火箭在自由空间由静止出发，火箭的排气速率为常量 $v_{0}$ ，求火箭的动量达到最大值时火箭的剩余质量。
\end{problem}
\clearpage


\begin{problem}{6.5.6}
一火箭无初速、垂直向上发射. 它靠喷射质量推进, 相对于火箭喷射速率 u 不变, 质量变化率为常量, 此常量由火箭的初加速度为零确定, 设重力加速度为常量. 求:

(1) 火箭的加速度与时间的关系；

(2) 火箭的高度与时间的关系(不必求出积分).
\end{problem}
\clearpage


\begin{problem}{6.5.7}
质量为 $M_0$ 并装了质量为 $m_0$ 的燃料的火箭船在均匀重力场中垂直起飞, 如图 6.33 所示. 相对于火箭船以速率 $u_0$ 喷射燃料, 在 $T_0$ 时间内全部喷出. 求 $T_0$ 时火箭船的速度.
\begin{center}
\begin{tikzpicture}[>=Stealth, thick]
  % Rocket body with rounded nose circle
  \draw (0,3.1) circle (0.24);
  \draw (-0.24,3.1) .. controls (-0.52,2.1) and (-0.72,1.1) .. (-0.72,0)
        -- (0.72,0)
        .. controls (0.72,1.1) and (0.52,2.1) .. (0.24,3.1);

  % Exhaust gas streams (broken downward arrows)
  \foreach \x in {-0.54,-0.27,0,0.27,0.54} {
    \draw (\x,-0.15) -- (\x,-0.70);
    \draw[->] (\x,-0.92) -- (\x,-1.55);
  }
  \node[right] at (0.62,-0.82) {$u_0$};

  % Gravity arrow
  \draw[->] (1.35,3.2) -- (1.35,2.05) node[midway, right] {$g$};
\end{tikzpicture}

\end{center}
\end{problem}
\clearpage


\begin{problem}{6.5.8}
一个瞬时质量为 $m$ 的火箭受到一个恒定推力 $F$ 的作用, 此推力是由发射推进剂取得的, 发射的相对速率较大, $\frac{\mathrm{d}m}{\mathrm{d}t}$ 较小, 推力总沿速度方向. 火箭沿在一个平面内的近似螺旋线的轨道从开始离地球中心距离 $r_1$ 运动到 $r_2, r_1$ 与地球半径 $r_0$ 接近, 那里的重力加速度为 $g_0, r_2 \gg r_0$ , 火箭离地球中心的角坐标为 $\varphi$ .

(1) 火箭的单位质量角动量是否是运动常数?

(2) 导出关于瞬时速率和重力加速度用 $r, r_{0}, \dot{r}$ 及 $\dot{\varphi}$ 表达的式子；

(3) 导出用上述各量及 $n$ 表示的 $\dot{r}$ 及 $\ddot{\varphi}$ 的表达式.
\end{problem}
\clearpage


\begin{problem}{6.5.9}
火箭在无引力场和无阻力的情况下做直线运动, 设火箭的初始质量为 $m_{0}$ , 最终质量为 $m_{1}$ , 喷射的相对速率为常量 $v_{r}$ . 求从开始到全部燃料烧完时推力所做的功.
\end{problem}
\clearpage


\begin{problem}{6.5.10}
一根其上端受到一个向上的力 $F$ 作用的线密度为 $\mu$ 的均质链条, 垂直地下放到桌面上, 如图6.34所示. 求桌面以上的链条端点高度 $h$ 所满足的运动微分方程.
\begin{center}
\begin{tikzpicture}[>=Stealth, thick]
  % Table surface
  \draw (-1.6,-0.35) -- (-1.6,0) -- (1.6,0) -- (1.6,-0.35);

  % Vertical chain links
  \foreach \i in {0,...,12} {
    \draw (0, 0.12 + \i*0.23) ellipse (0.075 and 0.115);
  }

  % Coiled chain links resting on the table
  \draw (-0.22,0.065) ellipse (0.13 and 0.065);
  \draw (-0.52,0.065) ellipse (0.14 and 0.065);
  \draw (-0.82,0.055) ellipse (0.12 and 0.055);
  \draw (0.23,0.065) ellipse (0.13 and 0.065);
  \draw (0.53,0.065) ellipse (0.14 and 0.065);
  \draw (0.81,0.055) ellipse (0.12 and 0.055);

  % Upward force F
  \draw[->] (0,3.05) -- (0,4.05) node[right] {$F$};

  % Height dimension h
  \draw[thin] (0.08,3.0) -- (1.75,3.0);
  \draw[<->, thin] (1.45,0) -- (1.45,3.0) node[midway, fill=white, inner sep=2pt] {$h$};
\end{tikzpicture}

\end{center}
\end{problem}
\clearpage


\begin{problem}{6.5.11}
火箭在真空和无引力情况下做直线运动,初始质量和最终质量分别为 $m_{0}$ 和 $m_{1}$ . 定义火箭的机械效率为燃料烧完后火箭的动能与所消耗的能量之比, 问齐奥尔可夫斯基数 $z = \frac{m_0}{m_1}$ 多大时火箭的效率最大. 设喷射速率 $v_{r}$ 不变.
\end{problem}
\clearpage


\begin{problem}{6.5.12}
物体沿水平轨道滑行, 燃气以不变的相对于物体的速率 $v_{r}$ 铅直地向下排出, 物体的初速度为 $v_{0}$ . 设物体的质量按规律 $m = m_{0} e^{-\alpha t}$ (其中 $\alpha$ 为常量) 变化, 滑动摩擦因数为 $\mu$ , 求此物体速度和位移的变化规律. 又问: 当 $\alpha$ 值多大时, 物体将以不变的速度运动?
\end{problem}
\clearpage


\begin{problem}{6.5.13}
一球状尘埃在均匀密度的水蒸气云雾中形成液滴,降落时液滴上的沉积率正比于单位时间内液滴扫过的体积.如果液滴在云雾中由静止开始运动,求液滴的加速度值.
\end{problem}
\clearpage


\begin{problem}{6.5.14}
一质量为 $m_0$ 、横截面积为 $A$ 的宇宙飞船以速度 $v_0$ 航行，穿过一密度为 $\rho$ 的处于静止状态的尘云，如尘埃粘贴到宇宙飞船上。求以后宇宙飞船的运动（设 $A$ 不随时间变化）。
\begin{center}
\begin{tikzpicture}[>=Stealth, thick, scale=0.82]
  % Cylindrical spaceship
  \draw (-3.2,0.5) .. controls (-3.45,0.22) and (-3.45,-0.22) .. (-3.2,-0.5);
  \draw (-3.2,0.5) -- (-1.35,0.5);
  \draw (-3.2,-0.5) -- (-1.35,-0.5);
  \draw (-1.35,0) ellipse (0.26 and 0.5);
  \node at (-1.35,0) {$A$};

  % Velocity arrow
  \draw[->] (-0.90,0) -- (-0.12,0) node[midway, above] {$v_0$};

  % Dust cloud particles (deterministic scatter avoiding text boxes)
  \foreach \x/\y in {
    0.35/1.15, 0.55/0.75, 0.42/0.25, 0.60/-0.22, 0.38/-0.78, 0.52/-1.18,
    0.85/1.22, 0.98/0.92, 0.78/0.48, 0.92/0.05, 0.80/-0.45, 0.95/-0.95,
    1.22/1.08, 1.15/0.62, 1.12/-0.12, 1.25/-0.68, 1.18/-1.22,
    1.52/1.25, 1.48/0.85, 1.55/-0.02, 1.45/-1.08,
    1.85/1.12, 1.92/0.82, 1.88/0.02, 1.78/-1.15,
    2.20/1.22, 2.28/0.88, 2.22/-0.05, 2.18/-0.98, 2.35/-1.22,
    2.55/1.05, 2.65/0.68, 2.52/0.22, 2.68/-0.25, 2.58/-0.72, 2.70/-1.12,
    2.95/1.20, 3.05/0.82, 2.88/0.42, 3.12/0.08, 2.95/-0.42, 3.08/-0.88,
    3.28/1.02, 3.35/0.52, 3.25/-0.15, 3.32/-0.65, 3.22/-1.18,
    0.68/-1.08, 2.42/0.04, 2.78/0.95, 2.85/-0.12, 1.68/-1.25
  } {
    \fill (\x,\y) circle (1.15pt);
  }

  % Cloud labels with clean white backing
  \node[fill=white, inner sep=2.5pt] at (1.85,0.42) {尘云};
  \node[fill=white, inner sep=2.5pt] at (1.85,-0.58) {密度 $\rho$};
\end{tikzpicture}

\end{center}
\end{problem}
\clearpage


\begin{problem}{6.5.15}
一股横截面积为 $A$ 、密度为 $\rho$ ，以绝对速度 $v_0$ 在水平方向运动的水流中，无弹性地射中一质量为 $m$ 的木块，即水离开木块时有相对于木块的水平速度分量为零。若木块与它在其上滑动的水平面间的摩擦因数为 $\mu$ 。求木块的最终速度。
\end{problem}
\clearpage


\begin{problem}{6.5.16}
一质量为 $M$ (无水时) 的水桶开始处于静止状态, 桶中装有质量为 $m$ 的水, 通过一根绳子施以恒力 $F$ 将桶从井中提上来, 桶中的水以不变的速率从桶中漏出来, 经过时间 $T$ 桶变成空的. 求变成空桶的瞬间桶的速度.
\end{problem}
\clearpage


\begin{problem}{6.5.17}
一个砂漏如图 6.36 所示, 置于一个磅秤上, 容器的质量为 $M$ , 砂子质量为 $m$ . 在 $t < 0$ 时, 砂子置于容器的上部; $t = 0$ 时, 打开中间的活门, 让砂子下落, 若单位时间内有 $\lambda$ (正的常数) 的砂子离开容器的上部, 自由下落的距离为 $h$ . 问 $t \geqslant 0$ 时, 磅秤的读数将如何变化? 设砂子刚到达容器底部时, 上部还有砂子.
\begin{center}
\begin{tikzpicture}[>=Stealth, thick]
  % Scale base (trapezoid)
  \draw (-1.05,-2.5) -- (1.05,-2.5) -- (0.62,-1.95);
  \draw (-1.05,-2.5) -- (-0.62,-1.95);

  % Scale dial
  \draw (0,-1.0) circle (1.15);
  \fill (0,-1.0) circle (2pt);
  \draw (0,-1.0) -- (34:0.82);

  % Hourglass container (lower and upper cones)
  \draw (-1.95,0.15) -- (1.95,0.15) -- (-1.10,2.25) -- (1.10,2.25) -- cycle;

  % Sand / falling stream dashed lines
  \draw[dashed] (-0.18,1.72) -- (0.18,1.72);
  \draw[dashed] (0,1.72) -- (0,0.34);
  \draw[dashed] (-0.72,0.15) .. controls (-0.25,0.38) and (0.25,0.38) .. (0.72,0.15);

  % Height h dimension
  \draw[dashed, thin] (0,1.40) -- (2.45,1.40);
  \draw[thin] (2.05,0.15) -- (2.45,0.15);
  \draw[<->, thin] (2.30,0.15) -- (2.30,1.40) node[midway, fill=white, inner sep=2pt] {$h$};
\end{tikzpicture}

\end{center}
\end{problem}
\clearpage


\section{位力定理}

\begin{problem}{6.6.1}
质点系中各质点的质量和矢径分别为 $m_{i}$ 和 $r_{i}$ ，分别受到净作用力 $F_{i}$ .

(1) 如果 $\sum_{i} m_{i} \dot{r}_{i} \cdot r_{i}$ 始终保持有限，证明位力定理：

\[
\overline {{{T}}} = - \frac {1}{2} \overline {{{\sum_ {i} F _ {i} \cdot r _ {i}}}}
\]

其中 T 为系统的总动能，“——”表示对时间求平均；

(2) 对于一个质点只受平方反比有心力的情况, 证明

\[
\overline {{T}} = - \frac {1}{2} \overline {{V}}
\]

其中 V 是势能.
\end{problem}
\clearpage


\begin{problem}{6.6.2}
用位力定理解 6.1.15 题

提示: 在质心平动参考系中对系统中一个质点使用位力定理.
\end{problem}
\clearpage


\begin{problem}{6.6.3}
有一群质量为 $m$ 的粒子, 它们围绕同一个中心做圆轨道运动, 每个粒子具有相同的动能 $T$ , 作用力只有相互间的万有引力. 求粒子数密度与离中心的距离 $r$ 间的函数关系 (假设密度是球对称的).
\end{problem}
\clearpage


\begin{problem}{6.6.4}
证明: 如果作用于质点系的力 $F_{i}$ 是非摩擦力 $F_{i}^{\prime}$ 和正比于速度的摩擦力 $f_{i}$ 之和, 位力只与 $F_{i}^{\prime}$ 有关, 没有来自 $f_{i}$ 的贡献.
\end{problem}
\clearpage

